\Needspace{11\baselineskip}
\section{Normalization for the Gauduchon equation}

By \cref{thm:stw-volume}, every smooth real function $F$ determines a scalar $b_\gamma(F)$, the constant $b$ of that theorem. As for the Monge--Amp\`ere equation, its variation is governed by the adjoint probability measure $\mu_u^\gamma$ of \cref{prop:gaud-volume-linearization} and by the concavity of the logarithmic determinant. To differentiate solutions, we fix a smooth positive volume form $\mathrm{d}V_0$ and use the normalization $\int_Xu\,\mathrm{d}V_0=0$ in place of $\sup_Xu=0$.

\begin{lemma}\label{lem:gaud-volume-smooth-dependence}
The solution $(u,b_\gamma(F))$, normalized by $\int_Xu\,\mathrm{d}V_0=0$, depends smoothly on the parameters of every smooth finite-dimensional family of functions $F$.
\end{lemma}
\begin{proof}
The map $\Gamma\mapsto\eta$ with $\eta^{n-1}=\Gamma$ is smooth on positive forms, since in a frame it is the inverse of the smooth map $\eta\mapsto\eta^{n-1}$, whose differential is invertible. The openness argument in the proof of \cref{thm:stw-volume}, with the $\mathrm{d}V_0$-mean in place of the $\gamma^n$-mean, gives a bounded inverse for the derivative with respect to the potential and the scalar between the relevant H\"older spaces. The implicit function theorem gives local smooth solution maps, and uniqueness shows that they agree on overlaps. Applying this for every order of H\"older regularity gives the assertion for smooth data.
\end{proof}

\Needspace{7\baselineskip}
For a smooth family of prescribed functions $F_t\in C^\infty(X,\mathbb R)$, with $t$ in an open interval of $\mathbb R$, let $u_t$ be the corresponding solutions with $\int_Xu_t\,\mathrm{d}V_0=0$, and set $Q_t=Q_{u_t}$ as in \eqref{eq:stw-cofactor-matrix-scaling}. The expression
\[
 \operatorname{tr}\bigl((Q_t^{-1}\dot Q_t)^2\bigr)
\]
is independent of the chosen frame and nonnegative, because $Q_t^{-1}\dot Q_t$ is similar to the Hermitian matrix $Q_t^{-1/2}\dot Q_t Q_t^{-1/2}$.

\Needspace{18\baselineskip}
\begin{theorem}\label{thm:gaud-volume-normalization-derivatives}
For a smooth family $F_t$ and its solutions,
\begin{align*}
 \dot b_\gamma(F_t)&=-\int_X\dot F_t\,\mathrm{d}\mu_{u_t}^\gamma,\\
 \ddot b_\gamma(F_t)&=-\int_X\ddot F_t\,\mathrm{d}\mu_{u_t}^\gamma
 -\frac1{n-1}\int_X
        \operatorname{tr}\bigl((Q_t^{-1}\dot Q_t)^2\bigr)
                          \,\mathrm{d}\mu_{u_t}^\gamma.
\end{align*}
In particular, $b_\gamma$ is concave, and along $F_t=F_0+tG$ its second derivative is strictly negative unless $G$ is constant. Moreover, for smooth real functions $F,G$ on $X$ and every real constant $a$,
\[
 b_\gamma(F+a)=b_\gamma(F)-a,
 \qquad |b_\gamma(G)-b_\gamma(F)|\leq\|G-F\|_\infty.
\]
\end{theorem}
\begin{proof}
By \eqref{eq:stw-cofactor-matrix-scaling},
\[
 \mathcal G(u)=\frac1{n-1}\log\det_\gamma Q_u.
\]
Differentiating $\mathcal G(u_t)=F_t+b_\gamma(F_t)$ once gives $\mathcal L_{u_t}\dot u_t=\dot F_t+\dot b_\gamma(F_t)$, and integration against the adjoint probability measure gives the first formula.

Because $\Gamma_u$ is affine in $u$, its second derivative along the family is $\mathcal D_\gamma\ddot u_t$. The formula for the second derivative of $\log\det$ used in the proof of \cref{thm:b-derivatives} therefore gives
\[
 \mathcal L_{u_t}\ddot u_t
 -\frac1{n-1}\operatorname{tr}\bigl((Q_t^{-1}\dot Q_t)^2\bigr)
 =\ddot F_t+\ddot b_\gamma(F_t).
\]
Integrating against $\mu_{u_t}^\gamma$ at this fixed parameter gives the second formula; no derivative of the measure is needed.

For affine $F_t$, we have $\ddot F_t=0$, and positivity of the density of $\mu_{u_t}^\gamma$ proves concavity. If the second derivative vanishes, then $\dot Q_t=0$, so $\dot\Gamma_{u_t}=0$ and $\mathcal L_{u_t}\dot u_t=0$. By \cref{prop:gaud-volume-linearization}, $\dot u_t$ is constant, and the first differentiated equation forces $G$ to be constant. Conversely, adding a constant to $F$ changes only $b$, by the opposite constant. Finally, integrate the first derivative along the segment from $F$ to $G$: since $\mu_{u_t}^\gamma$ is a probability measure, the derivative lies between $-\sup_X(G-F)$ and $-\inf_X(G-F)$, which proves the Lipschitz bound.
\end{proof}

\Needspace{6\baselineskip}
The Lipschitz bound extends $b_\gamma$ uniquely to continuous functions $F$. Constructing corresponding weak solutions would require, in addition, a volume measure for the $(n-1)$st roots of singular forms in the Aeppli ansatz, compactness of their potentials, and stability of the adjoint equation.

\Needspace{13\baselineskip}
On a surface with Gauduchon background form $\gamma_0$, conservation of the total mass makes the scalar and the adjoint probability measure explicit.

\Needspace{9\baselineskip}
\begin{example}\label{ex:gaud-volume-surface}
If $n=2$ and $\gamma_0$ is Gauduchon, then $\Gamma_u=\gamma_0+\ddc u=\eta_u$, and $\eta_u$ is also Gauduchon. The mass identity of \cref{thm:mass-invariance-characterization} gives $\int_X\eta_u^2=\int_X\gamma_0^2$. Therefore
\[
 b_\gamma(F)=\log\int_X\gamma_0^2-
                   \log\int_X\mathrm{e}^F\gamma^2,
 \qquad
 \mathrm{d}\mu_u^\gamma=\frac{\mathrm{e}^F\gamma^2}{\int_X\mathrm{e}^F\gamma^2}.
\]
The adjoint equation of \cref{prop:gaud-volume-form-adjoint} gives the same measure, because $\ddc\eta_u=0$ and hence $\rho_u$ is constant. Along $F+tG$, the second derivative of $b_\gamma$ is minus the variance of $G$ with respect to this probability measure.
\end{example}
